Read in Gauge Theory Gravity, the metric, Ricci tensor and covariant derivatives below are the derived metric, the contractions of the field strength (with $\Ric = \mathcal{R}$ as a map on physical vectors in the store's sign) and $D$ (Fact \claimref{synthesis.imported-gtg-dictionary}, clauses (i), (ii) and (iv)). \emph{Inputs.} From Premise \claimref{synthesis.c-loc-local-horizon-premise} only clause (Loc) (density $\eta$ per unit area, one universal constant) and clause (R$_{\mathrm{loc}}$) ($a/2\pi$ in the proper time of the boost observers) are used, together with its construction of $\mathcal{H}_{p,k}$ (vanishing expansion and shear at $p$). From Fact \claimref{synthesis.imported-jacobson-clausius-1995} we take the setting (local Rindler horizons through every point in every null direction, the approximate boost field $\chi$, the heat as boost energy of matter) and its hypothesis (c), which is (H1) here and which carries the focusing of the generators through the Raychaudhuri equation; we do not use the Fact's conclusion as a black box but rederive it from (H1) to (H4), so that every step is visible. The explicit equations are Fact \claimref{synthesis.imported-raychaudhuri-null-congruence}: for a twist-free congruence of null geodesics with tangent $k$ and affine parameter $\lambda$, the Raychaudhuri equation $d\theta/d\lambda = -\tfrac12\theta^{2} - |\sigma|^{2} - \Ric^{\circ}(k,k)$ with $|\sigma|^{2} \geq 0$ the squared shear, and the shear equation $d\sigma/d\lambda = -\theta\sigma - C(\cdot,k,\cdot,k)$ on the screen space, so that $\sigma(0) = 0$ gives $\sigma = O(\lambda)$ and $|\sigma|^{2} = O(\lambda^{2})$. The shear equation is the one step that is not inside the Fact. Both $\Ric^{\circ}(k,k)$ and $\theta$ are unchanged under $g \mapsto -g$, so they hold as written in signature $(1,3)$. The generators of the null hypersurface $\mathcal{H}_{p,k}$ are twist-free.

\emph{(1) Temperature.} In the flat tangent model at $p$ with coordinates $(t,x)$ in the plane of the boost, $\chi = \kappa(x\partial_{t} + t\partial_{x})$. Its orbit $x^{2} - t^{2} = \rho^{2}$, parametrised rationally by $t = \tfrac{\rho}{2}(u - u^{-1})$, $x = \tfrac{\rho}{2}(u + u^{-1})$, has $\dot t^{2} - \dot x^{2} = \rho^{2}/u^{2}$, unit velocity $U = \tfrac{u}{\rho}(\dot t, \dot x)$ and acceleration $\tfrac{u}{\rho}\dot U$ of square $-1/\rho^{2}$, so $a = 1/\rho$; and $g(\chi,\chi) = \kappa^{2}(x^{2} - t^{2}) = \kappa^{2}\rho^{2}$. The time of $\chi$ runs at $1/|\chi|$ times proper time, so energies and temperatures in it are $|\chi|$ times the proper ones: by (R$_{\mathrm{loc}}$), $T_{\chi} = (a/2\pi)|\chi| = (1/2\pi\rho)\,\kappa\rho = \kappa/2\pi$, the same on every orbit. On the past horizon $x = -t$ of the wedge, with $k = \partial_{t} - \partial_{x}$ and $\lambda = t < 0$, $\chi = \kappa(-\lambda\,\partial_{t} + \lambda\,\partial_{x}) = -\kappa\lambda k$, future-directed for $\lambda < 0$; for a field of nonvanishing field strength this holds up to $O(\lambda^{2})$ (H2), and curvature corrections to the flat tangent model are of higher order in the patch size, the order (H2) discards.

\emph{(2) Focusing.} Write $r(\lambda) = \Ric^{\circ}(k,k)$ and $t(\lambda) = T(k,k)$ along the generator through $p$, $r = r_{0} + r_{1}\lambda + \dots$, $t = t_{0} + t_{1}\lambda + \dots$ (Taylor expansions with remainders $O(\lambda^{2})$, since the metric is $C^{3}$ and $T$ is $C^{1}$), $|\sigma|^{2} = s_{2}\lambda^{2} + \dots$ (by the construction of $\mathcal{H}_{p,k}$ and the shear equation), and $\chi = -\kappa(\lambda + e_{2}\lambda^{2} + \dots)k$. The Raychaudhuri equation with $\theta(0) = 0$ gives $\theta = -r_{0}\lambda - \tfrac12 r_{1}\lambda^{2} + O(\lambda^{3})$, the terms $\theta^{2}$ and $|\sigma|^{2}$ first entering $\theta$ at $\lambda^{3}$. Per (H2) we work per unit cross-section area in a thin pencil about this generator, the cross-section shrunk to zero before $\lambda_{0}$ (neighbouring generators of a finite element have expansion $O(\epsilon)$ at $\lambda = 0$ for an element of size $\epsilon$). The cross-section area obeys $dA/d\lambda = \theta A$, so $A(\lambda)/A(0) = \exp\int_{0}^{\lambda}\theta = 1 - \tfrac12 r_{0}\lambda^{2} + O(\lambda^{3})$ and the area gained between $\lambda_{0}$ and $p$ is, per unit area at $p$, $\delta A = A(0) - A(\lambda_{0}) = \tfrac12 r_{0}\lambda_{0}^{2} + O(\lambda_{0}^{3})$. The heat is $\delta Q = \int_{\lambda_{0}}^{0} T(\chi,k)\,(A(\lambda)/A(0))\,d\lambda = -\kappa\int_{\lambda_{0}}^{0}(\lambda + e_{2}\lambda^{2})\,t(\lambda)\,(1 + O(\lambda^{2}))\,d\lambda = \tfrac12\kappa t_{0}\lambda_{0}^{2} + O(\lambda_{0}^{3})$. With $T_{\chi}\,\delta S = (\kappa/2\pi)\,\eta\,\delta A$ by (1), (Loc) and the differences reading of (H1),
\[
\delta Q - T_{\chi}\eta\,\delta A = \tfrac{\kappa}{2}\big(t_{0} - \tfrac{\eta}{2\pi}r_{0}\big)\lambda_{0}^{2} + O(\lambda_{0}^{3}),
\]
with no terms of order $\lambda_{0}^{0}$ or $\lambda_{0}^{1}$; $r_{1}, t_{1}, e_{2}$ first appear at $\lambda_{0}^{3}$ and $s_{2}$ at $\lambda_{0}^{4}$. Since (H1) holds at leading order for every small patch, the coefficient vanishes: $t_{0} = (\eta/2\pi)r_{0}$, and $\kappa \neq 0$ drops out. The conditions at $p$ are needed: if $|\sigma|^{2}(p) = s_{0} \neq 0$, the same computation gives $\theta = -(r_{0} + s_{0})\lambda + \dots$ and the coefficient $\tfrac{\kappa}{2}(t_{0} - \tfrac{\eta}{2\pi}(r_{0} + s_{0}))$; if $\theta(0) = \theta_{0} \neq 0$, then $\delta A = -\theta_{0}\lambda_{0} + O(\lambda_{0}^{2})$ while $\delta Q = O(\lambda_{0}^{2})$, a first-order defect $\kappa\eta\theta_{0}\lambda_{0}/2\pi$, so (H1) forces $\theta_{0} = 0$. As $k$ and $p$ are arbitrary, $\Ric^{\circ}(k,k) = (2\pi/\eta)T(k,k)$ for every null $k$ on $U$.

\emph{(3) Restatement in the store's language.} For null $k$, $\Sch(k,k) = \tfrac12\Ric(k,k) = -\tfrac12\Ric^{\circ}(k,k)$, so the own-ray class of Theorem \claimref{synthesis.einstein-condition-is-infinity-vacuum-incidence}(2) is $\omega_{k}(I_{g},k) = -\Sch(k,k) = \tfrac12 r_{0}$. With $4\pi G_{\mathrm{cl}} = \pi/\eta$, $\tfrac{\kappa}{2}(t_{0} - \tfrac{\eta}{2\pi}r_{0}) = \kappa\tfrac{\eta}{2\pi}(\tfrac{\pi}{\eta}t_{0} - \tfrac12 r_{0}) = \kappa\tfrac{\eta}{2\pi}(4\pi G_{\mathrm{cl}}\,T(k,k) - \omega_{k}(I_{g},k))$. This is an identity obtained by renaming $\tfrac12\Ric^{\circ}(k,k)$; it adds no content beyond the translation that theorem already records.

\emph{(4) Einstein's equations.} By (2), $(\Ric + 8\pi G_{\mathrm{cl}}T)(k,k) = -\Ric^{\circ}(k,k) + (2\pi/\eta)T(k,k) = 0$ for every real null $k$, which is the hypothesis $\omega_{k}(I_{g},k) = 4\pi G_{\mathrm{cl}}T(k,k)$ of Theorem \claimref{synthesis.einstein-condition-is-infinity-vacuum-incidence}(7) with $G = G_{\mathrm{cl}}$. That theorem gives $\Ric + 8\pi G_{\mathrm{cl}}T = (\Lambda + 4\pi G_{\mathrm{cl}}\tau)g$ with $\Lambda$ constant on the connected $U$: by polarisation (a symmetric form vanishing on the null cone is a multiple of $g$, in every dimension at least three, as the check confirms in dimensions three and four), then the trace and the contracted Bianchi identity with $\nabla^{a}T_{ab} = 0$ (H3), which needs $\Ric$ differentiable ($C^{3}$ metric) and $T$ in $C^{1}$. Four dimensions (H4) enter only here, through the trace: in dimension $n$ the same steps make $h - 8\pi G_{\mathrm{cl}}\tau/(n-2)$ constant, which is $h - 4\pi G_{\mathrm{cl}}\tau$ for $n = 4$ (and clause (7) is stated in four dimensions). In the other conventions: $\Ric^{\circ} = 8\pi G_{\mathrm{cl}}T - h g$ with $h = \Lambda + 4\pi G_{\mathrm{cl}}\tau$, so $R^{\circ} = 8\pi G_{\mathrm{cl}}\tau - 4h$ and $\Ric^{\circ} - \tfrac12R^{\circ}g - \Lambda g = 8\pi G_{\mathrm{cl}}T$; under $g \mapsto -g$, $\Ric^{\circ}$, $R^{\circ}g$ and $T_{ab}$ are unchanged and $\Lambda g$ changes sign, giving $G_{ab} + \Lambda g_{ab} = 8\pi G_{\mathrm{cl}}T_{ab}$ in signature $(3,1)$, Jacobson's form. Nothing in (C$_{\mathrm{loc}}$) or (H1) to (H4) constrains the constant of integration $\Lambda$.

\emph{(5) Newton's constant.} $8\pi G_{\mathrm{cl}} = 2\pi/\eta$ gives $G_{\mathrm{cl}} = 1/(4\eta)$, Jacobson's value. In the mode-count reading of version 1 of the premise, $\eta = 4\pi\nu\langle s\rangle/G$ per unit area, that is $4\pi\nu\langle s\rangle$ per Planck area, so $G_{\mathrm{cl}} = G/(16\pi\nu\langle s\rangle)$.

\emph{(6) Universality.} On the differences reading of (H1), let $\gamma = 2\pi/\eta$ be a function. Steps (2) and (4) up to polarisation are pointwise, so $\Ric + \gamma T = h g$ still holds. Then $\nabla^{a}\Ric_{ab} = \nabla_{b}h - \gamma\nabla^{a}T_{ab} - T_{ab}\nabla^{a}\gamma = \nabla_{b}h - T(\nabla\gamma)_{b}$; the trace gives $R = 4h - \gamma\tau$; and the contracted Bianchi identity $\nabla^{a}\Ric_{ab} = \tfrac12\nabla_{b}R$ (valid in either sign convention) gives $\nabla_{b}h - T(\nabla\gamma)_{b} = 2\nabla_{b}h - \tfrac12\gamma\nabla_{b}\tau - \tfrac12\tau\nabla_{b}\gamma$, that is $\nabla_{b}(h - \tfrac12\gamma\tau) = -T(\nabla\gamma)_{b}$. For constant $\eta$ this is the constancy of $\Lambda$ in (4); for dust $T = \rho\,u\otimes u$ it is $-\rho\,(u\cdot\nabla\gamma)\,u$, nonzero wherever $\gamma$ varies along the flow. On the integrated reading $S = \int\eta\,dA$, per unit cross-section $\delta S = \eta(p)A(0) - \eta(\lambda_{0})A(\lambda_{0}) = -(k\cdot\nabla\eta)\lambda_{0} + O(\lambda_{0}^{2})$ while $\delta Q = O(\lambda_{0}^{2})$, so the defect has the first-order term $(\kappa/2\pi)(k\cdot\nabla\eta)\lambda_{0}$; (H1) then forces $k\cdot\nabla\eta = 0$ for every null $k$, hence $\nabla\eta = 0$, already at step (2). The store takes the differences reading.

\emph{Remarks.} (a) The modular-flow clause of version 1 of the premise is not part of version 2 and is not used; it enters Jacobson's 2016 entanglement-equilibrium route (Fact \claimref{synthesis.imported-jacobson-entanglement-equilibrium-2016}), not this one. (b) $G_{\mathrm{cl}} = G$ is the condition $16\pi\nu\langle s\rangle = 1$, which is the area-law condition that defines the registered $\nu$ (Theorem \claimref{synthesis.clausius-route-newton-constant-is-residue}(iii)), when the horizon modes of the premise and the modes counted at the curvature cutoff are one set placed twice. If they were independent sets, the area-law calibration would halve $\eta$ and give $G_{\mathrm{cl}} = 2G$ (Theorem \claimref{synthesis.c-loc-is-a-second-placement-of-the-count}). Either way this is a consistency, not a determination of $G$: a value $\eta = 1/4$ taken from the Schwarzschild area law (Theorem \claimref{synthesis.c-loc-entropy-clause-curvature-independent-density}) already contains Einstein's $G$ through the Schwarzschild metric and temperature, and returns $G_{\mathrm{cl}} = G$. Theorem \claimref{synthesis.clausius-route-newton-constant-is-residue} is cited here as a reading, not used as a premise. (c) What selects Einstein's equations over their Ricci-polynomial extensions (which share the Killing-horizon data, Theorem \claimref{synthesis.black-hole-data-do-not-fix-field-equations}) is the assumption that the entropy density does not depend on curvature: a curvature-dependent density, as for a Wald entropy, is excluded by (Loc), and (6) shows where such a dependence would act. Only the leading order in the patch size is used, which is the order at which (C$_{\mathrm{loc}}$) is stated. (d) (H1) and (H3) are where the stance that gravity is not a quantised field enters: $T_{ab}$ is a given source and the heat is its boost energy; the derivation never quantises the metric.

\emph{The check.} Exact power series (GA.ps\_*) solve the Raychaudhuri equation and integrate the area and the heat for several rational values of $\kappa$, $\eta/2\pi$, the slopes $r_{1}, t_{1}$, the shear coefficient $s_{2}$ and the Killing error $e_{2}$: the defect has vanishing orders $0$ and $1$, its $\lambda_{0}^{2}$ coefficient is $\tfrac{\kappa}{2}(t_{0} - \tfrac{\eta}{2\pi}r_{0})$ in every case, the $\lambda_{0}^{3}$ coefficient varies with the nuisance values (control) and the shear alone moves only the $\lambda_{0}^{4}$ coefficient, and the zero set does not depend on $\kappa$; a shear or an expansion at $p$ shifts the defect as stated. On the integrated reading with $\eta/2\pi$ linear along the generator, the first-order defect is $\kappa$ times its slope, and it is absent on the differences reading. The temperature step is checked on the rational parametrisation of the Rindler orbit (acceleration squared $-1/\rho^{2}$, $|\chi|^{2} = \kappa^{2}\rho^{2}$, $a|\chi| = \kappa$) and on the past horizon. The null-null step is checked on eleven rational null vectors (rank nine, kernel $g$) and on seven in dimension three (rank five, kernel $g$), with the conversion to the store's form, to $\Ric^{\circ} - \tfrac12 R^{\circ}g - \Lambda g = 8\pi G_{\mathrm{cl}}T$ with $\Lambda = h - 4\pi G_{\mathrm{cl}}\tau$, and the restatement (3) on a generic rational $\Ric^{\circ}$ and $T$. The Bianchi step is solved pointwise in dimensions three to six (the coefficient $1/(n-2)$, equal to $\tfrac12$ only in four) and with a gradient of $\gamma$; $G_{\mathrm{cl}} = 1/(4\eta)$ and its mode-count reading are checked with $\pi$ and $\ln 2$ tracked as formal symbols.
